% ============================================================
%  The Particle Model — Lesson Booklet (one section per lesson: A Reading · B Theory · C Practice · D Above level)
%  Year 7 (Year 8 extension)  |  Victorian Curriculum F–10 Version 2.0
%  Compile with: xelatex three times, then python3 fill.py Particle_Model_Lesson_Booklet_Yr7
% ============================================================
\documentclass[11pt,a5paper,twoside]{article}
\newcommand{\wrcoverfile}{../../../covers/water-resources.jpg}\newcommand{\unitname}{Advanced Water and Resources}\newcommand{\unitcaps}{ADVANCED WATER AND RESOURCES}

% ── Fonts (compile with LuaLaTeX) ────────────────────────────────────────────
\usepackage{fontspec}
% Bundled fonts (in ./fonts/) so the document compiles anywhere with XeLaTeX/LuaLaTeX
\setmainfont{Lato}[                   % body — humanist sans (stands in for deck's Inter)
  Path=../fonts/, Extension=.ttf, Ligatures=TeX,
  UprightFont=Lato-Regular, ItalicFont=Lato-Italic,
  BoldFont=Lato-Bold, BoldItalicFont=Lato-BoldItalic ]
\newfontfamily\headingfont{Poppins}[  % display/headings — matches the slide decks
  Path=../fonts/, Extension=.ttf, Ligatures=TeX,
  UprightFont=Poppins-Regular, ItalicFont=Poppins-Italic,
  BoldFont=Poppins-Bold, BoldItalicFont=Poppins-BoldItalic ]
\setlength{\headheight}{22pt}

\usepackage[protrusion=true,expansion=false]{microtype}
\usepackage[a5paper, inner=17mm, outer=13mm, top=19mm, bottom=17mm, headheight=14pt, headsep=4mm, footskip=8mm]{geometry}
\usepackage{xcolor}
\usepackage{tikz}
\usetikzlibrary{arrows.meta, shapes.geometric, calc, positioning, decorations.pathreplacing}
\usepackage{tcolorbox}
\tcbuselibrary{skins, breakable, raster}
\usepackage{fancyhdr}
\usepackage{titlesec}
\usepackage{enumitem}
\usepackage{booktabs}
\usepackage{tabularx}
\usepackage{array}
\usepackage{colortbl}
\usepackage{multicol}
\usepackage{amsmath}
\usepackage{needspace}
% TikZ diagram macros are embedded below (just before \begin{document}) for easy editing.
\usepackage{longtable}
\usepackage{graphicx}
\usepackage{pifont}
\usepackage{lastpage}
\usepackage{calc}
\usepackage{mdframed}
\usepackage{setspace}
\usepackage{ragged2e}
% PDF bookmarks: one per lesson (made in \sectiondiv) plus the cover and
% contents, and links from the contents table to each lesson.
% implicit=false: hyperref adds no anchors of its own (an anchor on each
% list item nudged lists down 1.5 mm, which would undo the page fill).
% hidelinks keeps the links invisible, so the printed pages are unchanged.
\usepackage[implicit=false, hidelinks,
  pdftitle={Advanced Water and Resources},
  pdfsubject={Years 7–8 Science, Victorian Curriculum F–10 v2.0}]{hyperref}
\usepackage{bookmark} % writes the outline itself, which implicit=false switches off in hyperref
% hyperref 7.01 with implicit=false leaves \Hy@PageAnchorSlide undefined, which
% throws at \end{document} and stops fill.py (it builds with -halt-on-error).
\makeatletter\providecommand\Hy@PageAnchorSlide{}\makeatother

% ── Colour Palette ───────────────────────────────────────────────────────────
% Remapped to the slide-deck palette (names kept so body + diagrams cascade)
\definecolor{NavyBlue}   {HTML}{16305C}
\definecolor{MidBlue}    {HTML}{2E6FE0}
\definecolor{LightBlue}  {HTML}{E8F1FF}
\definecolor{TealGreen}  {HTML}{0F6E63}
\definecolor{LightGreen} {HTML}{E0F7F3}
\definecolor{GoldDark}   {HTML}{C77F12}
\definecolor{GoldLight}  {HTML}{FFF3DC}
\definecolor{PurpleDark} {HTML}{5B3FC0}
\definecolor{PurpleLight}{HTML}{EFE9FF}
\definecolor{OrangeDark} {HTML}{C9483D}
\definecolor{OrangeLight}{HTML}{FFE9E6}
\definecolor{MidGrey}    {HTML}{AEBED4}
\definecolor{LightGrey}  {HTML}{F2F5FA}
\definecolor{DarkGrey}   {HTML}{3A4759}
\definecolor{TableHead}  {HTML}{16305C}
\definecolor{RowAlt}     {HTML}{EEF4FF}
% Bright accent tints (deck) for bars, rules and chips
\definecolor{Blue}   {HTML}{2E6FE0}
\definecolor{Teal}   {HTML}{18A999}
\definecolor{Amber}  {HTML}{F5A623}
\definecolor{Coral}  {HTML}{FF6B5E}
\definecolor{Violet} {HTML}{7C5CE0}
\definecolor{Sky}    {HTML}{4DA3FF}
\definecolor{Ink}    {HTML}{16223A}
\definecolor{Line}   {HTML}{DDE5F0}
\definecolor{Paper}  {HTML}{FBFCFE}

% ── Page Style ───────────────────────────────────────────────────────────────
\pagestyle{fancy}
\fancyhf{}
\renewcommand{\headrulewidth}{0.4pt}
\renewcommand{\headrule}{\hbox to\headwidth{\color{NavyBlue}\leaders\hrule height \headrulewidth\hfill}}
\fancyhead[L]{\headingfont\small\color{NavyBlue}The Particle Model}
\fancyhead[R]{\headingfont\small\color{DarkGrey}Year 7 \textbar{} Victorian Curriculum F--10 v2.0}
\fancyfoot[C]{\headingfont\small\color{DarkGrey}\thepage\ \textbar{} \pageref{LastPage}}
\fancypagestyle{plain}{%
  \fancyhf{}
  \fancyfoot[C]{\headingfont\small\color{DarkGrey}\thepage\ \textbar{} \pageref{LastPage}}
  \renewcommand{\headrulewidth}{0pt}
}

% ── Section Formatting ───────────────────────────────────────────────────────
\titleformat{\section}
  {\headingfont\color{NavyBlue}\Large\bfseries}
  {}{0em}{}
  [{\color{Teal}\rule[0.6ex]{2.2em}{3pt}}{\color{Line}\rule[0.6ex]{\dimexpr\linewidth-2.2em\relax}{1pt}}]
\titlespacing{\section}{0pt}{18pt}{8pt}

\titleformat{\subsection}
  {\headingfont\color{MidBlue}\large\bfseries}
  {\thesubsection}{0.6em}{}
\titlespacing{\subsection}{0pt}{12pt}{4pt}

\titleformat{\subsubsection}
  {\headingfont\color{NavyBlue}\normalsize\bfseries}
  {}{0em}{}
\titlespacing{\subsubsection}{0pt}{8pt}{2pt}

% Deck-style eyebrow / kicker (optional use)
\newcommand{\kicker}[2][Teal]{%
  {\headingfont\bfseries\footnotesize\color{#1}%
   \textcolor{#1}{\rule[0.18ex]{1.6em}{2.2pt}}\hspace{0.5em}\MakeUppercase{#2}\par}%
  \vspace{2pt}}

% ── Custom tcolorbox Environments ────────────────────────────────────────────

% Section divider banner
\newtcolorbox{sectionbanner}[2]{
  enhanced, breakable=false,
  colback=NavyBlue, colframe=NavyBlue,
  fontupper=\color{white}\bfseries\large,
  left=4mm, right=4mm, top=3mm, bottom=3mm,
  title={\color{white}\Huge\bfseries #1},
  attach boxed title to top left={yshift=-2mm, xshift=4mm},
  boxed title style={colback=NavyBlue, colframe=NavyBlue},
  after upper={\par\smallskip\normalfont\normalsize\color{white!80!NavyBlue}\textit{#2}}
}

% Key info box (blue)
\newtcolorbox{infobox}[1]{
  breakable, arc=2.5mm, boxrule=0.5pt,
  colback=LightBlue, colframe=Blue!40,
  fonttitle=\headingfont\bfseries, coltitle=MidBlue,
  title={#1},
  left=4mm, right=3.5mm, top=2.5mm, bottom=2.5mm
}

% Scaffold box (yellow)
\newtcolorbox{scaffold}[1]{
  breakable, arc=2.5mm, boxrule=0.5pt,
  colback=GoldLight, colframe=Amber!40,
  fonttitle=\headingfont\bfseries, coltitle=GoldDark,
  title={#1},
  left=4mm, right=3.5mm, top=2.5mm, bottom=2.5mm
}

% Worked example box
\newtcolorbox{workedexbox}[1]{
  breakable, arc=2.5mm, boxrule=0.5pt,
  colback=white, colframe=Amber!50,
  fonttitle=\headingfont\bfseries, coltitle=GoldDark,
  title={#1},
  left=4mm, right=3.5mm, top=2.5mm, bottom=2.5mm
}
% Worked examples: \begin{workedex}{title} -> "Worked Example N.M --- title",
% auto-numbered per section and labelled we:N.M. Never type the number.
\newcounter{wenum}[section]
\renewcommand{\thewenum}{\thesection.\arabic{wenum}}
\newenvironment{workedex}[1]
  {\refstepcounter{wenum}\label{we:\thewenum}\begin{workedexbox}{Worked Example \thewenum\ --- #1}}
  {\end{workedexbox}}

% We do: straight after each worked example, a near-identical question the
% class works through together, with its own writing space. Numbered to
% match the worked example above it and labelled wd:N.M.
%   \begin{wedo}{title} question ... \answerbox{h} \end{wedo}
\newtcolorbox{wedobox}[1]{
  enhanced, breakable=false, arc=2.5mm, boxrule=0.5pt,
  colback=white, colframe=TealGreen!55,
  fonttitle=\headingfont\bfseries, coltitle=TealGreen!75!black,
  title={#1},
  left=4mm, right=3.5mm, top=2.5mm, bottom=2.5mm
}
\newenvironment{wedo}[1]
  {\par\needspace{9\baselineskip}\begin{wedobox}{We do \thewenum\ --- #1}\label{wd:\thewenum}\ignorespaces}
  {\end{wedobox}}

% Year 7 box (blue)
\newtcolorbox{yr7box}{
  breakable, arc=2.5mm, boxrule=0.5pt,
  colback=LightBlue, colframe=Blue!40,
  left=4mm, right=3.5mm, top=2.5mm, bottom=2.5mm
}

% Year 8 extension box (violet)
\newtcolorbox{yr8box}{
  breakable, arc=2.5mm, boxrule=0.5pt,
  colback=PurpleLight, colframe=Violet!40,
  left=4mm, right=3.5mm, top=2.5mm, bottom=2.5mm
}

% Answer lines box
\newcommand{\answerbox}[1]{%
  \begin{tcolorbox}[
    enhanced, breakable, arc=2mm,
    colback=Paper, colframe=Line,
    left=2mm, right=2mm, top=1mm, bottom=1mm,
    boxrule=1pt,
    height=#1
  ]
  \end{tcolorbox}%
}

% Experiment header
\newtcolorbox{expheader}[3]{
  breakable=false, arc=2.5mm,
  colback=RowAlt, colframe=Amber!50, boxrule=0.5pt,
  fonttitle=\headingfont\bfseries\large, coltitle=NavyBlue,
  title={Experiment #1: #2},
  left=4mm, right=3.5mm, top=2.5mm, bottom=2.5mm,
  after upper={\par\smallskip
    {\headingfont\bfseries\color{GoldDark}Aim:}\ \textit{#3}},
}

% Reading comprehension passage box (green)
\newtcolorbox{readingbox}[1]{
  breakable, arc=2.5mm, boxrule=0.5pt,
  colback=white, colframe=Teal!50,
  fonttitle=\headingfont\bfseries, coltitle=TealGreen,
  title={#1},
  left=4mm, right=3.5mm, top=2.5mm, bottom=2.5mm,
  fontupper=\normalsize
}

% Skill-builder box (for beginner scaffolding)
\newtcolorbox{skillbox}[1]{
  breakable, arc=2.5mm, boxrule=0.5pt,
  colback=RowAlt, colframe=Blue!40,
  fonttitle=\headingfont\bfseries, coltitle=MidBlue,
  title={#1},
  left=4mm, right=3.5mm, top=2.5mm, bottom=2.5mm
}

% ── Custom Commands ───────────────────────────────────────────────────────────

% Question: \question{marks}{text}  (marks can be empty)
% Numbered automatically as Q<section>.<n>, reset every section.
% Do NOT write numbers by hand.
\newcounter{qnum}[section]
\renewcommand{\theqnum}{\thesection.\arabic{qnum}}
\newcommand{\question}[2]{%
  \refstepcounter{qnum}%
  \needspace{3\baselineskip}%
  \medskip
  \noindent\label{q:\theqnum}\textbf{Q\theqnum}\ifx&#1&\else\ \textcolor{DarkGrey}{(#1 mark\ifnum#1=1\else s\fi)}\fi\quad #2%
  \medskip
}

% Inline marks note
\newcommand{\markscount}[1]{\hfill\textcolor{DarkGrey}{\textit{(#1 mark\ifnum#1=1\else s\fi)}}}

% ── Practice set counter and question macros ─────────────────────────────────
\newcounter{pqnum}
\newcommand{\practiceset}[1]{%
  \needspace{6\baselineskip}
  \setcounter{pqnum}{0}%
  \vspace{6pt}
  {\headingfont\bfseries\large\color{GoldDark} Practice Set --- #1}\\[-2pt]
  {\color{Amber}\rule{\linewidth}{1pt}}\\[2pt]
}
% \pq{question text}{working box height}
\newcommand{\pq}[2]{%
  \stepcounter{pqnum}%
  \needspace{4\baselineskip}
  \vspace{5pt}
  \noindent\textbf{\thepqnum.}\quad #1\par
  \vspace{2pt}
  \answerbox{#2}
}
% Two short questions side by side, each with its own box
\newcommand{\pqpair}[4]{%
  \needspace{4\baselineskip}
  \vspace{5pt}
  \noindent
  \begin{minipage}[t]{0.48\linewidth}
    \stepcounter{pqnum}\textbf{\thepqnum.}\quad #1\par\vspace{2pt}
    \begin{tcolorbox}[enhanced,arc=2mm,colback=Paper,colframe=Line,left=2mm,right=2mm,
      top=1mm,bottom=1mm,boxrule=1pt,height=#2]\end{tcolorbox}
  \end{minipage}\hfill
  \begin{minipage}[t]{0.48\linewidth}
    \stepcounter{pqnum}\textbf{\thepqnum.}\quad #3\par\vspace{2pt}
    \begin{tcolorbox}[enhanced,arc=2mm,colback=Paper,colframe=Line,left=2mm,right=2mm,
      top=1mm,bottom=1mm,boxrule=1pt,height=#4]\end{tcolorbox}
  \end{minipage}
}

% ── Comprehension question macro (compact, lined) ────────────────────────────
\newcounter{rcnum}
\newcommand{\rcreset}{\setcounter{rcnum}{0}}
\newcommand{\rcq}[2]{%
  \stepcounter{rcnum}%
  \needspace{3\baselineskip}
  \vspace{4pt}
  \noindent\textbf{C\thercnum.}\quad #1\par
  \vspace{2pt}
  \linedbox{#2}
}

% Section divider (big numbered banner): \sectiondiv{title}{summary}
%
% This IS the section: it steps the real section counter, so subsections,
% tasks, questions and figures inside it are numbered N.M automatically.
% Put \label{sec:N} straight after it. Never hand-type a section number.
\newcommand{\sectiondiv}[2]{%
  \needspace{5\baselineskip}%
  \refstepcounter{section}%
  % PDF bookmark for the lesson; its anchor, lesson.N.1, is what the contents table links to
  \pdfbookmark[1]{Lesson \thesection{} — #1}{lesson.\thesection}%
  \begin{tcolorbox}[
    enhanced, breakable=false, arc=3.5mm,
    colback=NavyBlue, colframe=NavyBlue, boxrule=0pt,
    sidebyside, sidebyside align=center,
    lefthand width=2.4cm,
    left=2mm, right=4mm, top=4mm, bottom=4mm
  ]
    \begin{center}
      {\headingfont\bfseries\color{white!35!NavyBlue}\fontsize{50}{50}\selectfont \thesection}
    \end{center}
    \tcblower
    {\headingfont\bfseries\color{white}\LARGE #1}\\[3pt]
    {\itshape\color{white!75!NavyBlue}\small #2}
  \end{tcolorbox}%
}

% Tasks: \task, \task[a] (part a of a NEW task), \taskpart{b} (next part
% of the same task), \exttask (Year 8 "Extension" — same counter).
% Numbered Task <section>.<n>, reset every section. Do NOT type numbers.
% Every task, question and figure also gets an automatic label
% (task:1.7b, q:2.4, fig:1.2, fig:E2.1) so page refs come from the .aux.
\newcounter{tasknum}[section]
\renewcommand{\thetasknum}{\thesection.\arabic{tasknum}}
\NewDocumentCommand{\task}{o}{\refstepcounter{tasknum}\label{task:\thetasknum\IfValueT{#1}{#1}}\textbf{Task \thetasknum\IfValueT{#1}{#1}}}
\newcommand{\taskpart}[1]{\label{task:\thetasknum#1}\textbf{Task \thetasknum#1}}
\newcommand{\exttask}{\refstepcounter{tasknum}\label{task:\thetasknum}\textbf{Extension \thetasknum}}

% Figure captions: \figcap[fig:label]{caption} -> "Figure N.M --- caption"
% \expfigcap{caption} -> "Figure EN.M" (experiment pages). Both auto-numbered.
\newcounter{fignum}[section]
\renewcommand{\thefignum}{\thesection.\arabic{fignum}}
\newcounter{expfignum}[section]
\renewcommand{\theexpfignum}{E\thesection.\arabic{expfignum}}
\newcommand{\figcap}[2][]{%
  \begin{center}
    \refstepcounter{fignum}\label{fig:\thefignum}\if\relax\detokenize{#1}\relax\else\label{#1}\fi
    \small\itshape\color{DarkGrey} Figure \thefignum\ --- #2
  \end{center}%
}
\newcommand{\expfigcap}[1]{%
  \begin{center}
    \refstepcounter{expfignum}\label{fig:\theexpfignum}%
    \small\itshape\color{DarkGrey} Figure \theexpfignum\ --- #1
  \end{center}%
}

% Diagram include
\newcommand{\diagramfig}[3]{%  {filename}{width fraction}{caption}
  \begin{center}
    \includegraphics[width=#2\linewidth]{#1}\\[4pt]
    \small\itshape\color{DarkGrey} #3
  \end{center}%
}

% Table column types
\newcolumntype{L}[1]{>{\RaggedRight\arraybackslash}p{#1}}
\newcolumntype{C}[1]{>{\Centering\arraybackslash}p{#1}}

% Table header colour
\newcommand{\thead}[1]{\cellcolor{TableHead}\color{white}\bfseries #1}
\newcommand{\talt}{\rowcolor{RowAlt}}
\newcommand{\tmid}{\rowcolor{MidGrey}}

% Inline year tags
\newcommand{\tagyr}[1]{%
  \tikz[baseline=-0.5ex]\node[
    rounded corners=2pt,
    fill=#1!20, draw=#1, text=#1,
    inner xsep=4pt, inner ysep=1pt,
    font=\scriptsize\bfseries
  ]{Yr #1 only};}

% Draw answer lines
\newcommand{\anslines}[1]{%
  \vspace{2pt}%
  \foreach \n in {1,...,#1}{%
    \par\noindent\rule{\linewidth}{0.3pt}\vspace{3pt}%
  }%
}

% ── Document ─────────────────────────────────────────────────────────────────
% ── Optional-lesson note, shown under the lesson banner ──────────────────────
\newcommand{\optionalnote}[1]{%
  % the title is a bold lead-in on the first line rather than a title bar,
  % which keeps Lesson 11's reading on the lesson's first page
  \begin{tcolorbox}[enhanced, breakable=false, arc=2.5mm, boxrule=0.5pt,
    colback=PurpleLight, colframe=Violet!50,
    left=4mm, right=3.5mm, top=2mm, bottom=2mm]
  {\headingfont\bfseries\color{PurpleDark}Optional lesson \textbullet{} Year 8 extension.}\enspace #1
  \end{tcolorbox}\vspace{2pt}}


% ══════════════════════════════════════════════════════════════════════════════
%  LESSON STRUCTURE (lesson booklet, Oct 2026)
%  Every lesson runs in the same four parts, in this order:
%    \lessonpart{A}{Reading and comprehension}{Teal}   \label{read:N}
%    \lessonpart{B}{Theory and examples}{Blue}         \label{theory:N}
%    \lessonpart{C}{Practice}{Amber}                    \label{prac:N}
%    \lessonpart{D}{Above level}{Violet}                \label{ext:N}
%  Above-level / higher-order questions live ONLY in Part D, at the end.
% ══════════════════════════════════════════════════════════════════════════════
\newcommand{\lessonpart}[3]{%
  \par\if#1A\needspace{10\baselineskip}\vspace{10pt}\else
    \if#1C\notesfill\fi\clearpage\fi
  \begin{tcolorbox}[enhanced, breakable=false, arc=2mm, boxrule=0pt,
    colback=#3!14, borderline west={4pt}{0pt}{#3},
    left=4mm, right=3mm, top=1.6mm, bottom=1.6mm]
  {\headingfont\bfseries\color{#3!60!black}\large Part #1 \textbullet{} #2}
  \end{tcolorbox}\vspace{2pt}}

% Key vocabulary at the start of each lesson: students copy each definition
% from the vocabulary slide in the lesson's starter. First row = red scaffold.
\newcommand{\vocabhead}{%
  \par\vspace{6pt}%
  \begin{tcolorbox}[enhanced, breakable=false, arc=2mm, boxrule=0pt,
    colback=NavyBlue!10, borderline west={4pt}{0pt}{NavyBlue},
    left=4mm, right=3mm, top=1.6mm, bottom=1.6mm]
  {\headingfont\bfseries\color{NavyBlue}\large Key vocabulary}
  \end{tcolorbox}\vspace{2pt}%
  \noindent Copy the definition of each word from the slide. The first one has been done for you.\par}

% Curriculum codes under a lesson banner
\newcommand{\lessoncodes}[1]{%
  {\small\headingfont\color{DarkGrey}Victorian Curriculum 2.0:\ \textbf{#1}}\par\vspace{4pt}}

% Two short tasks side by side, each with its own working box
% \taskpair{text}{box height}{text}{box height}
% The two boxes line up: the shorter question is padded by the difference in
% text height (measured in a scratch box, task counter restored, no label).
\newsavebox\tpbox
\newlength\tpA \newlength\tpB
\newcount\tpsaved
\newcommand\tpmeasure[2]{%
  \tpsaved=\value{tasknum}%
  \sbox\tpbox{\vbox{\hsize=0.485\linewidth \linewidth=\hsize \parskip=0pt
    \renewcommand\label[1]{}\task #1\par}}%
  \setlength{#2}{\dimexpr\ht\tpbox+\dp\tpbox\relax}%
  \setcounter{tasknum}{\tpsaved}}
\newcommand{\taskpair}[4]{%
  \par\needspace{9\baselineskip}\vspace{4pt}%
  \tpmeasure{#1}{\tpA}\tpmeasure{#3}{\tpB}%
  \noindent
  \edef\abgrp{\number\numexpr\value{abox}+1\relax}%
  \begin{minipage}[t]{0.485\linewidth}
    \task #1\par\ifdim\tpB>\tpA\vspace{\dimexpr\tpB-\tpA\relax}\fi\vspace{2pt}\answerbox{#2}
  \end{minipage}\hfill
  \begin{minipage}[t]{0.485\linewidth}
    \task #3\par\ifdim\tpA>\tpB\vspace{\dimexpr\tpA-\tpB\relax}\fi\vspace{2pt}\answerbox{#4}
  \end{minipage}\par\gdef\abgrp{0}\vspace{4pt}}
% Two parts of one question side by side, with their answer boxes level:
% \ptpair{a}{part (a) text}{box height}{b}{part (b) text}{box height}.
% The shorter part's text gets space below it to match the longer one.
\newcommand\ptmeasure[3]{%
  \sbox\tpbox{\vbox{\hsize=0.485\linewidth \linewidth=\hsize \parskip=0pt \pt{#1}#2\par}}%
  \setlength{#3}{\dimexpr\ht\tpbox+\dp\tpbox\relax}}
\newcommand{\ptpair}[6]{%
  \ptmeasure{#1}{#2}{\tpA}\ptmeasure{#4}{#5}{\tpB}%
  \par\noindent
  \edef\abgrp{\number\numexpr\value{abox}+1\relax}%
  \begin{minipage}[t]{0.485\linewidth}
    \pt{#1}#2\par\ifdim\tpB>\tpA\vspace{\dimexpr\tpB-\tpA\relax}\fi\answerbox{#3}
  \end{minipage}\hfill
  \begin{minipage}[t]{0.485\linewidth}
    \pt{#4}#5\par\ifdim\tpA>\tpB\vspace{\dimexpr\tpA-\tpB\relax}\fi\answerbox{#6}
  \end{minipage}\par\gdef\abgrp{0}}

% Answer boxes never split across a page (a half box on its own page is useless)
% ── Page filling ────────────────────────────────────────────────────────────
% Every writing box gets an id. Its page is logged to \jobname.abp; fill.py
% measures the empty space left at the foot of each page and shares it out
% among that page's boxes via \jobname.abx (\abset{id}{extra}). Side-by-side
% boxes share a group so they grow together.
\newcounter{abox}
\newwrite\abpos \immediate\openout\abpos=\jobname.abp
\newcommand{\abset}[2]{\expandafter\gdef\csname abx@#1\endcsname{#2}}
\InputIfFileExists{\jobname.abx}{}{}
\def\abgrp{0}
\newcommand{\abnext}{\stepcounter{abox}%
  \edef\abcur{\ifcsname abx@\theabox\endcsname\csname abx@\theabox\endcsname\else 0pt\fi}}
\newcommand{\ablog}{\edef\abtmp{\noexpand\write\abpos{\theabox\space\abgrp\space\noexpand\thepage}}\abtmp}
\newcommand{\abh}[1]{\dimexpr#1+\abcur\relax}
\renewcommand{\answerbox}[1]{%
  \abnext
  \begin{tcolorbox}[enhanced, breakable=false, arc=2mm, colback=Paper, colframe=Line,
    left=2mm, right=2mm, top=1mm, bottom=1mm, boxrule=1pt, height=\abh{#1}]
  \end{tcolorbox}\ablog%
}
% \linedbox{height}: an answer box ruled in dotted thirds. Lines are \thirdrow/3 apart and run dotted, dotted,
% solid (the baseline), repeating, and end on one dotted line. The number of
% rows is worked out from the box's final height, so the ruling still fills the
% box after fill.py has grown it; the rows are centred in the box.
\newlength\thirdrow \setlength\thirdrow{9mm}
\tikzset{thirdsolid/.style={draw=DarkGrey!70, line width=0.5pt},
         thirddot/.style={draw=DarkGrey!70, line width=0.6pt, dash pattern=on 0.9pt off 1.6pt}}
\newlength\lbH \newlength\lbO \newcount\lbN
\newcommand{\linedbox}[1]{%
  \abnext
  % interior height = box height less the 1pt frame top and bottom
  \setlength\lbH{\dimexpr\abh{#1}-2pt\relax}%
  % whole rows that fit with at least 3 mm spare top and bottom (at least one)
  \lbN=\numexpr(\dimexpr\lbH-6mm-\thirdrow/2\relax)/\thirdrow\relax
  \ifnum\lbN<1 \lbN=1 \fi
  \setlength\lbO{\dimexpr(\lbH-\lbN\thirdrow)/2\relax}%
  \begin{tcolorbox}[enhanced, breakable=false, arc=2mm, colback=Paper, colframe=Line,
    left=2mm, right=2mm, top=1mm, bottom=1mm, boxrule=1pt, height=\abh{#1},
    underlay={\begin{tcbclipinterior}
      \foreach \i in {1,...,\the\lbN}{
        \draw[thirddot] ([xshift=3mm,yshift=-\lbO-(\i-1)*\thirdrow]interior.north west) -- ([xshift=-3mm,yshift=-\lbO-(\i-1)*\thirdrow]interior.north east);
        \draw[thirddot] ([xshift=3mm,yshift=-\lbO-(\i-1)*\thirdrow-\thirdrow/3]interior.north west) -- ([xshift=-3mm,yshift=-\lbO-(\i-1)*\thirdrow-\thirdrow/3]interior.north east);
        \draw[thirdsolid] ([xshift=3mm,yshift=-\lbO-(\i-1)*\thirdrow-2*\thirdrow/3]interior.north west) -- ([xshift=-3mm,yshift=-\lbO-(\i-1)*\thirdrow-2*\thirdrow/3]interior.north east);}
      \draw[thirddot] ([xshift=3mm,yshift=-\lbO-\lbN\thirdrow]interior.north west) -- ([xshift=-3mm,yshift=-\lbO-\lbN\thirdrow]interior.north east);
    \end{tcbclipinterior}}]
  \end{tcolorbox}\ablog%
}
% Every written-response box is ruled the same way (Part A comprehension, We do,
% Part C tasks and Part D extensions, including side-by-side pairs). Space for
% mathematical working uses \workbox{height} instead and stays plain, as do
% \drawbox and the write-in tables.
\let\workbox\answerbox   % plain box: mathematical working (calculations) stays unruled
\let\answerbox\linedbox
% \drawbox{height}{contents}: white box for a diagram or drawing
\newcommand{\drawbox}[2]{%
  \abnext
  \begin{tcolorbox}[colback=white, colframe=MidGrey, boxrule=0.5pt, height=\abh{#1}]#2\end{tcolorbox}\ablog}
% Notes box at the end of Part B: fills whatever is left of the page (if worth it)
\newcommand{\notesfill}{\par
  \ifdim\dimexpr\pagegoal-\pagetotal\relax>3.2cm
    \vspace*{\fill}%
    \begin{tcolorbox}[enhanced, breakable=false, arc=2mm, colback=Paper, colframe=Line,
      left=2mm, right=2mm, top=1mm, bottom=1mm, boxrule=1pt,
      height=\dimexpr\pagegoal-\pagetotal-1.0cm\relax]
      {\small\headingfont\color{DarkGrey}Notes}
    \end{tcolorbox}%
  \fi}



% Centred table: wrap every content table in \begin{ctab} ... \end{ctab}
% (figures already centre themselves inside their \diag... macros)
% ctab also registers write-in tables (S column) for page filling: the extra
% height is per row, logged with the row count.
\newenvironment{ctab}{\par\smallskip\begingroup\centering
  \abnext\let\tbcur\abcur\global\wtrows0 }{\par
  \ifnum\wtrows>0 \edef\abtmp{\noexpand\write\abpos{\theabox\space\abgrp\space\noexpand\thepage\space\the\wtrows}}\abtmp\fi
  \endgroup\smallskip}

% ── Booklet formatting rules (Oct 2026) ─────────────────────────────────────
% 1. No marks are shown anywhere in the booklet.
\renewcommand{\markscount}[1]{}
% 2. A task's heading sits on its own line; its text starts on the next
%    line, left aligned. Parts (a), (b)… each start a new line: use \pt{a}.
\RenewDocumentCommand{\task}{o}{\par\noindent\refstepcounter{tasknum}\label{task:\thetasknum\IfValueT{#1}{#1}}\textbf{Task \thetasknum\IfValueT{#1}{#1}}\par\nobreak\noindent\ignorespaces}
\renewcommand{\taskpart}[1]{\par\noindent\label{task:\thetasknum#1}\textbf{Task \thetasknum#1}\par\nobreak\noindent\ignorespaces}
\renewcommand{\exttask}{\par\noindent\refstepcounter{tasknum}\label{task:\thetasknum}\textbf{Extension \thetasknum}\par\nobreak\noindent\ignorespaces}
\newcommand{\pt}[1]{\par\noindent\textbf{(#1)}\ \ignorespaces}
% Workstation pages (Experiment 1): numbered method steps with a hanging
% indent, and a one-line safety reminder under the materials list.
\newcommand{\wsstep}[1]{\par\noindent\hangindent=1.6em\hangafter=1\makebox[1.6em][l]{\textbf{#1}}\ignorespaces}
\newcommand{\wssafety}[1]{\par\smallskip{\small\textcolor{ScaffoldRed}{\headingfont\bfseries Safety:}\ #1}\par}
% 3. Tables have vertical rules between columns and down both sides
%    (write | between and around the columns, e.g. {|L{3cm}|C{2cm}|}).
%    Keep the normal padding at the outer edges (no @{} there): colortbl paints
%    header colour one \tabcolsep past each edge, so trimming the padding makes
%    the header overhang the grid.
\arrayrulecolor{MidGrey}
\setlength{\arrayrulewidth}{0.6pt}
% 4. Write-in tables: W/Z are left/centred columns that centre text
%    vertically; put the S column first (see below).
\newcolumntype{W}[1]{>{\RaggedRight\arraybackslash}m{#1}}
\newcolumntype{Z}[1]{>{\Centering\arraybackslash}m{#1}}
\def\tbcur{0pt}
\newcount\wtrows
\newcommand{\wh}{\global\advance\wtrows1 \rule[\dimexpr-0.75cm-\tbcur/2\relax]{0pt}{\dimexpr1.5cm+\tbcur\relax}}
% S = zero-width first column that makes every row of a write-in table
% at least 1.5 cm tall with its text vertically centred. Start each data
% row with "&"; start the header row with "\hd &" so it stays slim.
\newif\ifnowh
\newcolumntype{S}{@{}>{}c<{\ifnowh\else\wh\fi\global\nowhfalse}@{\hspace{\tabcolsep}}}
\newcommand{\hd}{\global\nowhtrue}
% 5. The first row of a multi-row write-in table is a worked scaffold in red.
\definecolor{ScaffoldRed}{HTML}{C0392B}
\newcommand{\scaf}[1]{\textcolor{ScaffoldRed}{\itshape #1}}
% 6. Tables are a light grid (rules between rows and columns) on white —
%    row shading is off, because colortbl paints shading over column rules.
\renewcommand{\talt}{}
\renewcommand{\tmid}{}
% ══════════════════════════════════════════════════════════════════════════════
%  DIAGRAM DEFINITIONS — The Particle Model
%  Self-contained TikZ macros, no external image files. Each \diag... macro is
%  called from the body by name.
% ══════════════════════════════════════════════════════════════════════════════
\usetikzlibrary{arrows.meta,calc,positioning,patterns,backgrounds}
\tikzset{
  thin arrow/.style={-{Stealth[length=6pt,width=4pt]}, line width=1.4pt},
}

% ── Lesson 1: the same three jars holding a solid, a liquid and a gas ───────────
% The solid keeps its own shape, the liquid takes the jar's shape with a flat top,
% the gas fills the whole jar. No particles yet: those arrive in Lesson 4.
\newcommand{\jar}[1]{% jar outline with its base at (#1,0)
  \draw[line width=1.2pt, Ink] (#1-1.25,2.9) -- (#1-1.25,0.15) arc (180:270:0.15)
    -- (#1+1.1,0) arc (270:360:0.15) -- (#1+1.25,2.9);}
\newcommand{\diagStates}{%
\begin{center}
\begin{tikzpicture}[font=\small, x=1cm, y=1cm]
  % solid: a block sitting in the corner
  \fill[GoldLight] (-0.95,0.02) rectangle (0.15,1.25);
  \draw[GoldDark, line width=1.1pt] (-0.95,0.02) rectangle (0.15,1.25);
  \jar{0}
  \node[font=\bfseries, text=NavyBlue] at (0,-0.45) {Solid};
  \node[align=center, font=\footnotesize, text=DarkGrey] at (0,-1.05) {keeps its own shape\\and volume};
  % liquid: takes the jar's shape, flat top
  \begin{scope}[xshift=4.2cm]
    \fill[Sky!35] (-1.25,0.15) arc (180:270:0.15) -- (1.1,0) arc (270:360:0.15) -- (1.25,1.1) -- (-1.25,1.1) -- cycle;
    \draw[MidBlue, line width=1.1pt] (-1.25,1.1) -- (1.25,1.1);
    \jar{0}
    \node[font=\bfseries, text=NavyBlue] at (0,-0.45) {Liquid};
    \node[align=center, font=\footnotesize, text=DarkGrey] at (0,-1.05) {takes the jar's shape,\\keeps its volume};
  \end{scope}
  % gas: fills the whole (sealed) jar
  \begin{scope}[xshift=8.4cm]
    \fill[Violet!16] (-1.25,0.15) arc (180:270:0.15) -- (1.1,0) arc (270:360:0.15) -- (1.25,2.9) -- (-1.25,2.9) -- cycle;
    \draw[line width=2.2pt, DarkGrey] (-1.4,2.9) -- (1.4,2.9);
    \jar{0}
    \node[font=\scriptsize, text=DarkGrey, anchor=south] at (0,2.95) {lid};
    \node[font=\bfseries, text=NavyBlue] at (0,-0.45) {Gas};
    \node[align=center, font=\footnotesize, text=DarkGrey] at (0,-1.05) {spreads out to fill\\the whole jar};
  \end{scope}
\end{tikzpicture}
\end{center}}

% ── Lesson 1, Worked Example 1.2 and We do 1.2: the same syringe, before and after
\newcommand{\diagSyringe}[3]{% {start mL}{end mL}{contents: gas|liquid}
\begin{center}\resizebox{\linewidth}{!}{%
\begin{tikzpicture}[font=\scriptsize, x=0.11cm, y=1cm]
  \foreach \row/\vol/\lab in {0/#1/before, -1.15/#2/after}{
    \begin{scope}[yshift=\row cm]
      \ifx#3g\fill[Violet!16] (0,0) rectangle (\vol,0.7);\else\fill[Sky!35] (0,0) rectangle (\vol,0.7);\fi
      \draw[line width=1pt, Ink] (0,0) rectangle (60,0.7);
      \draw[line width=1pt, Ink] (0,0.25) -- (-4,0.25) (0,0.45) -- (-4,0.45);
      \fill[Ink] (-5.5,0.15) rectangle (-4,0.55);
      \foreach \t in {0,10,...,60}{\draw[Ink] (\t,0.7) -- (\t,0.85); \node[anchor=south, font=\tiny] at (\t,0.82) {\t};}
      \fill[DarkGrey] (\vol,0.04) rectangle (\vol+1.6,0.66);
      \draw[line width=2pt, DarkGrey] (\vol+1.6,0.35) -- (\vol+26,0.35);
      \draw[line width=2pt, DarkGrey] (\vol+26,0.05) -- (\vol+26,0.65);
      \node[anchor=east, font=\scriptsize\bfseries, text=DarkGrey] at (-7,0.35) {\lab};
    \end{scope}}
  \node[font=\tiny, text=DarkGrey, anchor=west] at (61,0.35) {mL};
  \node[font=\tiny, text=DarkGrey, anchor=north] at (-4.75,-1.2) {sealed tip};
\end{tikzpicture}}
\end{center}}

% needspace.sty reserves room by planting a penalty -100 with 9 lines of
% temporary stretch. On a page that can't be filled tightly that fake break
% always wins, so a We do that fitted was pushed over and left half a page
% empty. This version breaks only when the room really isn't there, and
% ragged bottom lets every other page run on until something won't fit;
% fill.py then grows the writing boxes into what is left at the foot.
\renewcommand{\needspace}[1]{\par\ifdim\dimexpr\pagegoal-\pagetotal\relax<\dimexpr#1\relax\newpage\fi}
\raggedbottom
% Part D sits in a breakable yr8box, where \needspace can't see the page (the box
% is built off the page first). Wrap each extension in keepext so it is one
% unsplittable block: the yr8box then breaks between extensions, never inside one.
\newenvironment{keepext}{\par\begin{tcolorbox}[blankest, breakable=false]}{\end{tcolorbox}\par}
\usetikzlibrary{patterns}
\providecommand{\lessonart}{\thesection}

\providecommand{\unitname}{The Particle Model}\providecommand{\unitcaps}{THE PARTICLE MODEL}
% ══════════════════════════════════════════════════════════════════════════════
%  2027 WORKBOOK STYLE — A5 TOPIC BOOKLET (added by make_a5.py)
%  Same look as the Space A5 trial and the A4 experiment booklets: colour
%  cover, black-and-white inside, Outfit headings, a lesson opener band with
%  the lesson's ink symbol, boxes told apart by frame and icon:
%    I do ● solid · We do ●● dashed · You do ●●● heavy · key idea / reading ◆
%    thin · Above level ▲ double · Country ◎ dotted
%  The body still names the master's colours; they are mapped to ink here.
% ══════════════════════════════════════════════════════════════════════════════
\usepackage{pdflscape}
\usetikzlibrary{patterns}

% A picture wider than the A5 text block is scaled down to fit it.
\newsavebox\fitb
\newenvironment{fitbox}{\begin{lrbox}{\fitb}}{\end{lrbox}%
  \ifdim\wd\fitb>\linewidth\resizebox{\linewidth}{!}{\usebox\fitb}\else\usebox\fitb\fi}

% ── Fonts: Outfit for headings ──────────────────────────────────────────────
\renewfontfamily\headingfont{Outfit}[
  Path=../fonts/, Extension=.ttf, Ligatures=TeX,
  UprightFont=Outfit-Regular, BoldFont=Outfit-SemiBold,
  ItalicFont=Outfit-Regular, ItalicFeatures={FakeSlant=0.18},
  BoldItalicFont=Outfit-SemiBold, BoldItalicFeatures={FakeSlant=0.18} ]

% ── Black and white inside ──────────────────────────────────────────────────
\definecolor{Ink}{HTML}{111111}
\colorlet{NavyBlue}{Ink}\colorlet{MidBlue}{Ink}\colorlet{TealGreen}{Ink}
\colorlet{GoldDark}{Ink}\colorlet{PurpleDark}{Ink}\colorlet{OrangeDark}{Ink}
\colorlet{TableHead}{Ink}
\colorlet{Blue}{Ink}\colorlet{Teal}{Ink}\colorlet{Amber}{Ink}\colorlet{Coral}{Ink}
\colorlet{Violet}{Ink}\colorlet{Sky}{Ink}
\colorlet{LightBlue}{white}\colorlet{LightGreen}{white}\colorlet{GoldLight}{white}
\colorlet{PurpleLight}{white}\colorlet{OrangeLight}{white}\colorlet{RowAlt}{white}
\colorlet{LightGrey}{white}\colorlet{Paper}{white}
\colorlet{DarkGrey}{black!72}\colorlet{MidGrey}{black!45}\colorlet{Line}{black!45}
\arrayrulecolor{black!60}
% Particle kinds stay distinguishable without colour: white, light grey,
% dotted and dark.
\tikzset{
  pA/.style={circle, fill=white, draw=Ink, line width=0.6pt, minimum size=3.2mm, inner sep=0pt},
  pB/.style={circle, fill=black!35, draw=Ink, line width=0.6pt, minimum size=3.2mm, inner sep=0pt},
  pC/.style={circle, preaction={fill=white}, pattern=dots, draw=Ink, line width=0.6pt, minimum size=3.2mm, inner sep=0pt},
  pD/.style={circle, fill=black!75, draw=Ink, line width=0.6pt, minimum size=3.2mm, inner sep=0pt}}

% ── Page style: mirrored, running head on the outside ───────────────────────
\fancypagestyle{fancy}{\fancyhf{}%
  \renewcommand{\headrulewidth}{0.4pt}%
  \renewcommand{\headrule}{\hbox to\headwidth{\color{Ink}\leaders\hrule height \headrulewidth\hfill}}%
  \fancyhead[LE]{\headingfont\footnotesize \unitname}%
  \fancyhead[RO]{\headingfont\footnotesize\leftmark}%
  \fancyfoot[LE,RO]{\headingfont\footnotesize\thepage}}
\fancypagestyle{plain}{\fancyhf{}\renewcommand{\headrulewidth}{0pt}%
  \fancyfoot[LE,RO]{\headingfont\footnotesize\thepage}}
\fancypagestyle{opener}{\fancyhf{}\renewcommand{\headrulewidth}{0pt}%
  \fancyfoot[LE,RO]{\headingfont\footnotesize\thepage}}
\pagestyle{fancy}\markboth{}{}
\raggedbottom
\setlength{\emergencystretch}{2em}

\titleformat{\subsection}{\headingfont\large\bfseries}{\thesubsection}{0.6em}{}
\titleformat{\subsubsection}{\headingfont\normalsize\bfseries}{}{0em}{}

% ── Lesson opener band (replaces the navy section banner) ────────────────────
\renewcommand{\sectiondiv}[2]{%
  \clearpage\thispagestyle{opener}%
  \refstepcounter{section}\markboth{\thesection\ \textbullet\ #1}{\thesection\ \textbullet\ #1}%
  \pdfbookmark[1]{Lesson \thesection{} — #1}{lesson.\thesection}\hypertarget{lesson.\thesection.1}{}%
  \noindent\begin{minipage}[b]{\dimexpr\linewidth-36mm\relax}
    {\headingfont\bfseries\scriptsize\addfontfeatures{LetterSpace=14}\unitcaps\ \textbullet\ TOPIC}\par\vspace{1mm}
    {\headingfont\bfseries\fontsize{40}{40}\selectfont\thesection}\par\vspace{1.5mm}
    {\headingfont\bfseries\fontsize{15}{18}\selectfont\RaggedRight #1\par}\vspace{1mm}
    {\footnotesize\itshape\RaggedRight #2\par}
  \end{minipage}\hfill
  \begin{minipage}[b]{34mm}\includegraphics[width=34mm]{../art/l\lessonart-ink.pdf}\end{minipage}\par
  \vspace{2.5mm}\noindent\rule{\linewidth}{1.6pt}\par\vspace{1mm}}
\renewcommand{\lessoncodes}[1]{{\small\itshape VC 2.0: #1}\par\vspace{2pt}}

% Lesson parts: a heading and a rule; B, C and D start a new page.
\renewcommand{\lessonpart}[3]{%
  \par\if#1A\needspace{10\baselineskip}\vspace{8pt}\else
    \if#1C\notesfill\fi\clearpage\fi
  {\headingfont\bfseries\Large Part #1 \textbullet\ #2}\par\vspace{0.5mm}\noindent\rule{\linewidth}{0.6pt}\par\vspace{1mm}}
\renewcommand{\vocabhead}{%
  \par\vspace{4pt}{\headingfont\bfseries\Large Key vocabulary}\par\vspace{1pt}%
  \noindent{\small Copy the meaning of each word from the slide. The first one has been done for you.}\par}
\renewcommand{\thead}[1]{\headingfont\bfseries #1}

% ── Box system ──────────────────────────────────────────────────────────────
\newcommand{\icodots}[1]{\tikz[baseline=-0.6ex]{\foreach \k in {1,...,#1}{\fill[white] (\k*1.9mm,0) circle (0.7mm);}}}
\newcommand{\icokey}{\tikz[baseline=-0.6ex]{\fill[white] (0,0.9mm)--(0.9mm,0)--(0,-0.9mm)--(-0.9mm,0)--cycle;}}
\newcommand{\icoup}{\tikz[baseline=-0.6ex]{\fill[white] (-1mm,-0.8mm)--(1mm,-0.8mm)--(0,1mm)--cycle;}}
\newcommand{\icocountry}{\tikz[baseline=-0.6ex]{\draw[white,line width=0.5pt] (0,0) circle (1mm); \fill[white] (0,0) circle (0.45mm);}}
\newcommand{\icoread}{\tikz[baseline=-0.6ex]{\draw[white,line width=0.6pt] (-1.1mm,-0.8mm) rectangle (1.1mm,0.8mm); \draw[white,line width=0.5pt] (0,-0.8mm) -- (0,0.8mm);}}
\newcommand{\tagicon}[1]{\tikz[baseline=-0.6ex]\node[fill=Ink, rounded corners=0.8mm, inner xsep=1.2mm, inner ysep=1mm]{#1};}
\tcbset{wbbase/.style={enhanced, colback=white, arc=1.2mm,
    left=3mm, right=3mm, top=4.2mm, bottom=2.4mm, before skip=10pt, after skip=7pt,
    attach boxed title to top left={xshift=3mm, yshift=-2.4mm},
    boxed title style={colback=Ink, colframe=Ink, arc=0.8mm, boxrule=0pt,
                       left=1.6mm, right=1.8mm, top=0.6mm, bottom=0.6mm},
    fonttitle=\headingfont\bfseries\footnotesize, coltitle=white}}
\renewtcolorbox{workedexbox}[1]{wbbase, breakable, colframe=Ink, boxrule=1pt,
  title={\icodots{1}\ I do \textbullet\ #1}}
\renewenvironment{workedex}[1]
  {\refstepcounter{wenum}\label{we:\thewenum}\begin{workedexbox}{Worked example \thewenum\ --- #1}}
  {\end{workedexbox}}
\renewtcolorbox{wedobox}[1]{wbbase, breakable=false, frame hidden,
  borderline={1.1pt}{0pt}{Ink, dash pattern=on 4pt off 2.6pt},
  title={\icodots{2}\ #1}}
\renewenvironment{wedo}[1]
  {\par\begin{wedobox}{We do \thewenum\ --- #1}\label{wd:\thewenum}\ignorespaces}
  {\end{wedobox}}
\renewtcolorbox{yr7box}{wbbase, breakable, colframe=Ink, boxrule=2.4pt, top=5mm,
  title={\icodots{3}\ You do \textbullet\ Check your understanding}}
\renewtcolorbox{yr8box}{wbbase, breakable, frame hidden,
  borderline={0.6pt}{0pt}{Ink}, borderline={0.6pt}{1.6pt}{Ink},
  title={\icoup\ Above level \textbullet\ Year 8 extension}}
\renewtcolorbox{infobox}[1]{wbbase, breakable, colframe=Ink, boxrule=0.5pt, title={\icokey\ #1}}
\renewtcolorbox{scaffold}[1]{wbbase, breakable, colframe=Ink, boxrule=0.5pt, title={\icokey\ #1}}
\renewtcolorbox{skillbox}[1]{wbbase, breakable, colframe=Ink, boxrule=0.5pt, title={\icokey\ #1}}
\renewtcolorbox{readingbox}[1]{wbbase, breakable, colframe=Ink, boxrule=0.5pt, title={\icoread\ #1}}
\renewtcolorbox{expheader}[3]{wbbase, breakable=false, colframe=Ink, boxrule=1pt,
  title={\icokey\ Experiment #1: #2},
  after upper={\par\smallskip{\headingfont\bfseries Aim:}\ \textit{#3}}}
\newtcolorbox{countrybox}[1]{wbbase, breakable, frame hidden,
  borderline={1.4pt}{0pt}{Ink, line cap=round, dash pattern=on 0pt off 2.4pt},
  title={\icocountry\ #1}}

% ── Padding to a multiple of 4 (copier booklet mode) ────────────────────────
\newcommand{\notespage}{\clearpage\markboth{Notes}{Notes}%
  {\headingfont\bfseries\Large Notes}\par\vspace{3mm}%
  \noindent\begin{tikzpicture}\foreach \y in {0,...,17}{\draw[black!35, line width=0.4pt] (0,-\y*9mm) -- (\linewidth,-\y*9mm);}\end{tikzpicture}}
\newcount\padleft
\newcommand{\padtofour}{\clearpage
  \pgfmathtruncatemacro\pagesused{\the\value{page}-1}%
  \pgfmathtruncatemacro\padn{mod(4-mod(\pagesused,4),4)}%
  \padleft=\padn\relax
  \loop\ifnum\padleft>0 \notespage\advance\padleft by -1 \repeat
  \clearpage}

% ══════════════════════════════════════════════════════════════════════════════
%  WATER CYCLE AND RESOURCES — COVER
%  The unit's cover art (covers/water-resources.jpg), full bleed. The A5 lesson
%  booklet sets \wrcoverfile one folder deeper.
% ══════════════════════════════════════════════════════════════════════════════
\providecommand{\wrcoverfile}{../../covers/water-resources.jpg}
\newcommand{\wrcover}{%
\begin{tikzpicture}[remember picture, overlay]
  \node[anchor=north] at (current page.north) {\includegraphics[width=\paperwidth]{\wrcoverfile}};
\end{tikzpicture}}
% ── A5: every question full page width (stacked, never side by side) ───────
\renewcommand{\taskpair}[4]{%
  \par\needspace{6\baselineskip}\vspace{2pt}%
  \task #1\par\vspace{2pt}\answerbox{#2}%
  \par\needspace{6\baselineskip}%
  \task #3\par\vspace{2pt}\answerbox{#4}\par\vspace{2pt}}
\renewcommand{\ptpair}[6]{%
  \pt{#1}#2\par\answerbox{#3}%
  \pt{#4}#5\par\answerbox{#6}}
\renewcommand{\pqpair}[4]{%
  \needspace{4\baselineskip}\vspace{5pt}%
  \stepcounter{pqnum}\noindent\textbf{\thepqnum.}\quad #1\par\vspace{2pt}
  \begin{tcolorbox}[enhanced,arc=2mm,colback=Paper,colframe=Line,left=2mm,right=2mm,
    top=1mm,bottom=1mm,boxrule=1pt,height=#2]\end{tcolorbox}
  \stepcounter{pqnum}\noindent\textbf{\thepqnum.}\quad #3\par\vspace{2pt}
  \begin{tcolorbox}[enhanced,arc=2mm,colback=Paper,colframe=Line,left=2mm,right=2mm,
    top=1mm,bottom=1mm,boxrule=1pt,height=#4]\end{tcolorbox}}

\begin{document}
\pdfbookmark[1]{Cover}{cover}
\setlength{\parindent}{0pt}
\setlength{\parskip}{4pt}

% ---- Cover: the unit's cover art, full bleed, in colour ------------------
\thispagestyle{empty}
\wrcover
\begin{tikzpicture}[remember picture, overlay]
  
  \node[anchor=north east, fill=white, draw=Ink, line width=0.6pt, rounded corners=1.2mm,
        inner xsep=2mm, inner ysep=1.8mm, text width=30mm, font=\headingfont\footnotesize]
    at ([xshift=-4mm, yshift=-4mm]current page.north east)
    {\textbf{Name}\par\vspace{4mm}\hrule\vspace{2mm}\textbf{Class}\par\vspace{4mm}\hrule};
  \node[anchor=south east, fill=white, draw=Ink, line width=0.6pt, rounded corners=1.2mm,
        inner xsep=2.6mm, inner ysep=2mm, font=\headingfont\bfseries, align=right]
    at ([xshift=-4mm, yshift=6mm]current page.south east)
    {Advanced Water and Resources\\[-1pt]{\normalfont\headingfont\scriptsize Topics 1--4}};
\end{tikzpicture}
\mbox{}
\clearpage

% ---- Inside cover ---------------------------------------------------------
\thispagestyle{plain}
{\headingfont\bfseries\fontsize{22}{26}\selectfont Advanced Water and Resources}\par\vspace{1mm}
{\headingfont\large Topics 1--4 \textbullet\ Above the Year 7 level}\par\vspace{1mm}
{\small\itshape VC 2.0: VC2S8U09 \textbullet\ VC2S10U11 \textbullet\ VC2S10U15 \textbullet\ VC2S10U16}\par\vspace{4mm}
\textbf{Teacher}\ \hrulefill\par

\begin{countrybox}{Acknowledgement of Country}
This booklet was created on the unceded Country of the Wurundjeri Woi-wurrung people of the Kulin Nation. We acknowledge Elders past, present and emerging.
\end{countrybox}

\begin{infobox}{How this booklet works}
\textbf{Advanced Water and Resources} goes deeper into the ideas of the lesson books: more powerful models, the reasoning behind them, and harder problems that bring several ideas together. Each topic builds on the lessons named under its title. Each opens with its key vocabulary, then:\par\smallskip
\renewcommand{\arraystretch}{1.25}
\begin{tabular}{@{}p{24mm}p{70mm}@{}}
\textbf{Part A} & a reading with comprehension questions \\
\textbf{Part B} & the deeper idea, worked through \\
\textbf{Part C} & practice on your own \\
\textbf{Part D} & a challenge that pulls the ideas together \\
\end{tabular}

\smallskip
\begin{tabular}{@{}p{24mm}p{70mm}@{}}
\textbf{I do} \tagicon{\icodots{1}} & a worked example (solid frame) \\
\textbf{We do} \tagicon{\icodots{2}} & you try one with help (dashed frame) \\
\textbf{Reading} \tagicon{\icoread} & a passage to read (thin frame) \\
\end{tabular}
\end{infobox}

\clearpage
\markboth{Contents}{Contents}
{\headingfont\bfseries\LARGE Contents}\par\vspace{1mm}\noindent\rule{\linewidth}{0.6pt}\par\vspace{2mm}
\renewcommand{\arraystretch}{1.45}
\begin{tabular}{|C{1.1cm}|L{5.2cm}|C{1.7cm}|C{1.4cm}|}
  \hline\thead{Topic} & \thead{Title} & \thead{Deepens lesson} & \thead{Page} \\ \hline
  1 & \hyperlink{lesson.1.1}{Energy, the Water Cycle and Climate} & 1 & \hyperlink{lesson.1.1}{p.\ \pageref{sec:1}} \\ \hline
  2 & \hyperlink{lesson.2.1}{Resources Across Their Life Cycle} & 3 & \hyperlink{lesson.2.1}{p.\ \pageref{sec:2}} \\ \hline
  3 & \hyperlink{lesson.3.1}{Water Footprints} & 4 & \hyperlink{lesson.3.1}{p.\ \pageref{sec:3}} \\ \hline
  4 & \hyperlink{lesson.4.1}{Energy Efficiency and the Grid} & 5 & \hyperlink{lesson.4.1}{p.\ \pageref{sec:4}} \\ \hline
  --- & \hyperlink{glossary}{Glossary} & & \hyperlink{glossary}{\pageref{sec:glossary}} \\ \hline
\end{tabular}
\renewcommand{\arraystretch}{1}
\clearpage
\setcounter{section}{0}\renewcommand{\lessonart}{1}
\sectiondiv{Energy, the Water Cycle and Climate}{The Sun drives the cycle \textbullet{} Energy in evaporation and condensation \textbullet{} How long water stays \textbullet{} A warmer world}\label{sec:1}
\lessoncodes{Deepens Lesson 1 \textbullet{} VC2S8U09 \textbullet{} VC2S10U11}

\vocabhead
\begin{ctab}
\begin{tabular}{|SW{3.4cm}|W{7.0cm}|}
  \hline\hd & \thead{Word} & \thead{Definition (copy from the slide)} \\ \hline
  & \scaf{\textbf{Store}\label{voc:store}} & \scaf{A place where water is held for a time: the ocean, ice, groundwater, the air.} \\ \hline
  & \textbf{Flux}\label{voc:flux} & \\ \hline
  & \textbf{Residence time}\label{voc:residence-time} & \\ \hline
  & \textbf{Transpiration}\label{voc:transpiration} & \\ \hline
  & \textbf{Groundwater}\label{voc:groundwater} & \\ \hline
\end{tabular}
\end{ctab}

\lessonpart{A}{Reading and comprehension}{Teal}\label{read:1}

\begin{readingbox}{Reading --- The Sun's Water Pump}
Every year the Sun lifts about half a million cubic kilometres of water from the oceans, lakes, soil and leaves into the air. Evaporation takes energy: each gram of water needs about 2400\,J to escape as vapour, which is why sweating cools you. When the vapour condenses into cloud droplets, that energy is released again, warming the air around it. A thunderstorm is powered by the energy released as its clouds condense.

Scientists describe the water cycle as \textbf{stores}, where water is held, and \textbf{fluxes}, the flows between them. Plants are a big flux: water drawn up through their roots and evaporated from their leaves is \textbf{transpiration}. A single large tree can transpire hundreds of litres on a hot day.

How long water stays in a store is its \textbf{residence time}. A water particle stays in the air for about nine days before it falls as rain. In the deep ocean it may stay for thousands of years. Some \textbf{groundwater}, held in rock deep under inland Australia, fell as rain hundreds of thousands of years ago.

Warmer air can hold more water vapour, about 7\% more for each degree Celsius. As the climate warms, more water evaporates, and when it rains, more can fall at once. Scientists expect both heavier downpours and longer droughts in many places.
\end{readingbox}

\rcreset
\rcq{Where does the energy to evaporate water come from, and where does it go when the vapour condenses?}{3cm}
\rcq{What is the difference between a store and a flux? Give an example of each.}{2.5cm}
\rcq{Why can deep groundwater be called ``fossil water''?}{2.5cm}
\rcq{Why might a warmer climate bring both heavier rain and worse droughts?}{3cm}

\lessonpart{B}{Theory and examples}{Blue}\label{theory:1}

\subsection{Energy in the Water Cycle}\label{sub:a1.1}

\begin{workedex}{A canvas water bag}
Before refrigerators, people hung canvas bags of water on cars and verandahs. Water seeping through the canvas evaporated from the outside, and the water inside stayed cool on a hot day. Explain.

\smallskip
To evaporate, the water particles on the outside of the bag need energy. They take it from the bag and the water inside, so the water inside loses energy and its temperature falls. The hotter and drier the air, the faster the evaporation and the stronger the cooling.
\end{workedex}

\begin{wedo}{Thunderstorms}
Explain, using energy, why a large cloud forming on a summer afternoon can grow into a violent thunderstorm.
\answerbox{4.5cm}
\end{wedo}

\subsection{Residence Time}\label{sub:a1.2}

\begin{skillbox}{Residence time}
\[ \text{residence time} = \frac{\text{amount in the store}}{\text{rate of flow out of the store}} \]
\end{skillbox}

\begin{workedex}{A farm dam}
A farm dam holds 6\,million litres. Cattle, evaporation and leakage remove 40\,000\,L a day, and no rain falls. How long will it last?

\smallskip
$6\,000\,000 \div 40\,000 = \mathbf{150\ days}$: about five months without rain.
\end{workedex}

\begin{wedo}{Water in the air}
The air holds about 13\,000\,km\textsuperscript{3} of water, and about 1400\,km\textsuperscript{3} falls as rain and snow every day. Calculate the residence time of water in the air. Then explain why pumping groundwater that took 100\,000 years to collect is a different problem from using river water.
\workbox{5.5cm}
\end{wedo}

\lessonpart{C}{Practice}{Amber}\label{prac:1}

\begin{keepext}
\task Clearing a large forest can reduce the rainfall in the area nearby. Use the water cycle to suggest why.
\answerbox{5cm}
\end{keepext}
\begin{keepext}
\task Sea water is salty, but rain is fresh. Explain how the water cycle separates the water from the salt, and name the separation technique that copies it.
\answerbox{5cm}
\end{keepext}


\lessonpart{D}{Challenge}{Violet}\label{ext:1}

\task Most of Melbourne's water comes from rain falling on forested catchments that flows into reservoirs. The climate in these catchments is expected to become warmer and a little drier. Explain, as a chain of causes, how this could reduce the water reaching the reservoirs by more than the drop in rainfall alone. Use the words \emph{evaporation}, \emph{transpiration}, \emph{soil} and \emph{runoff}.
\answerbox{8cm}

\clearpage

% ══════════════════════════════════════════════════════════════════════════════
%  TOPIC 2 — RESOURCES ACROSS THEIR LIFE CYCLE               (deepens Lesson 3)
% ══════════════════════════════════════════════════════════════════════════════
\setcounter{section}{1}\renewcommand{\lessonart}{3}
\sectiondiv{Resources Across Their Life Cycle}{Cradle to grave \textbullet{} Embodied energy \textbullet{} Comparing choices with numbers \textbullet{} The circular economy}\label{sec:2}
\lessoncodes{Deepens Lesson 3 \textbullet{} VC2S8U09 \textbullet{} VC2S10H03}

\vocabhead
\begin{ctab}
\begin{tabular}{|SW{3.4cm}|W{7.0cm}|}
  \hline\hd & \thead{Word} & \thead{Definition (copy from the slide)} \\ \hline
  & \scaf{\textbf{Life cycle}\label{voc:life-cycle}} & \scaf{Every stage of a product's life: getting the raw materials, making it, using it and getting rid of it.} \\ \hline
  & \textbf{Life cycle assessment}\label{voc:life-cycle-assessment} & \\ \hline
  & \textbf{Embodied energy}\label{voc:embodied-energy} & \\ \hline
  & \textbf{Circular economy}\label{voc:circular-economy} & \\ \hline
  & \textbf{E-waste}\label{voc:e-waste} & \\ \hline
\end{tabular}
\end{ctab}

\lessonpart{A}{Reading and comprehension}{Teal}\label{read:2}

\begin{readingbox}{Reading --- The Hidden Story of a Can}
An aluminium drink can is in your hand for a few minutes. Its \textbf{life cycle} is far longer. Bauxite ore is mined, often in Queensland or Western Australia, refined into alumina, and then smelted into aluminium using huge amounts of electricity. The metal is rolled, shaped, printed, filled and trucked to a shop. After you drink it, the can is buried, burnt, or recycled.

The total energy used across all these stages is the can's \textbf{embodied energy}. Recycling aluminium saves about 95\% of the energy needed to make new metal from ore, because the hardest step, smelting the ore, is skipped. A \textbf{life cycle assessment} adds up the energy, water, land and pollution of every stage, so that two products, or two choices, can be compared fairly.

In a \textbf{circular economy}, products are designed from the start to be repaired, reused and finally recycled, so materials go round and round instead of ending in landfill. \textbf{E-waste}, old phones, computers and batteries, is a big test: it holds valuable gold, copper and rare metals, but is hard to take apart.
\end{readingbox}

\rcreset
\rcq{List the stages in the life cycle of an aluminium can.}{2.5cm}
\rcq{Why does recycling aluminium save so much energy?}{2.5cm}
\rcq{What does a life cycle assessment add up, and why?}{2.5cm}
\rcq{How is a circular economy different from ``make, use, throw away''?}{2.5cm}

\lessonpart{B}{Theory and examples}{Blue}\label{theory:2}

\subsection{Comparing Choices with Numbers}\label{sub:a2.1}

\begin{workedex}{Glass or plastic?}
A refillable glass bottle takes 5\,MJ (megajoules) to make and 0.2\,MJ to wash each time it is reused. A single-use plastic bottle takes 2\,MJ to make. Compare them over 20 uses.

\smallskip
Glass: $5 + 20 \times 0.2 = 5 + 4 = 9$\,MJ. \quad Plastic: $20 \times 2 = 40$\,MJ.\\
Over 20 uses the glass bottle uses \textbf{less than a quarter} of the energy. Over one use, plastic wins.
\end{workedex}

\begin{wedo}{The keep cup}
A ceramic cup takes 14\,MJ to make and 0.1\,MJ to wash each time. A paper cup takes 0.5\,MJ. After how many uses does the ceramic cup become the better choice? (Hint: find when $14 + 0.1n = 0.5n$.)
\workbox{5cm}
\end{wedo}

\subsection{What a Life Cycle Assessment Leaves Out}\label{sub:a2.2}

\begin{workedex}{Energy is not everything}
A life cycle assessment says paper bags use more energy to make than plastic bags. Does that settle which is better?

\smallskip
No. Energy is only one measure. Plastic bags last for centuries and harm wildlife if they escape into the sea; paper breaks down in months. Paper uses more water and land to grow trees. A fair comparison looks at several impacts, and the answer can depend on how many times each bag is used.
\end{workedex}

\begin{wedo}{Cotton T-shirts}
Growing the cotton for one T-shirt can use about 2700\,L of water. Describe two other stages of a T-shirt's life cycle and an impact of each. Suggest one change that would most reduce its total impact.
\answerbox{5cm}
\end{wedo}

\lessonpart{C}{Practice}{Amber}\label{prac:2}

\begin{keepext}
\task Recycling an aluminium can uses far less energy than making a new one from its ore. Explain why recycling metals reduces both energy use and the risks of extraction.
\answerbox{5cm}
\end{keepext}
\begin{keepext}
\task Electric cars use no petrol, but their batteries need mined lithium. Weigh up whether electric cars reduce the risks of extracting resources.
\answerbox{5cm}
\end{keepext}


\lessonpart{D}{Challenge}{Violet}\label{ext:2}

\task Australians replace their phones about every two to three years, and most old phones end up in drawers or landfill. Should the law require phone makers to take back every old phone and recover its materials? Make a reasoned case, considering the resources in a phone, the energy and risks of mining them, the cost, and who should pay.
\answerbox{9cm}

\clearpage

% ══════════════════════════════════════════════════════════════════════════════
%  TOPIC 3 — WATER FOOTPRINTS                                (deepens Lesson 4)
% ══════════════════════════════════════════════════════════════════════════════
\setcounter{section}{2}\renewcommand{\lessonart}{4}
\sectiondiv{Water Footprints}{Measuring a household's water \textbullet{} Virtual water in food \textbullet{} Planning to use less}\label{sec:3}
\lessoncodes{Deepens Lesson 4 \textbullet{} VC2S8U09 \textbullet{} VC2S10I04}

\vocabhead
\begin{ctab}
\begin{tabular}{|SW{3.4cm}|W{7.0cm}|}
  \hline\hd & \thead{Word} & \thead{Definition (copy from the slide)} \\ \hline
  & \scaf{\textbf{Water footprint}\label{voc:water-footprint}} & \scaf{The total amount of fresh water used to produce the things a person, family or country uses.} \\ \hline
  & \textbf{Virtual water}\label{voc:virtual-water} & \\ \hline
  & \textbf{Flow rate}\label{voc:flow-rate} & \\ \hline
  & \textbf{Per capita}\label{voc:per-capita} & \\ \hline
  & \textbf{Target}\label{voc:target} & \\ \hline
\end{tabular}
\end{ctab}

\lessonpart{A}{Reading and comprehension}{Teal}\label{read:3}

\begin{readingbox}{Reading --- The Water You Eat}
During the Millennium Drought, Melburnians were asked to use no more than 155 litres of water each per day at home. That \textbf{target}, Target 155, counted showers, toilets, washing machines and gardens: the water you can see coming out of a tap.

But far more water is used to make the things we buy. The water used to grow, process and transport a product is its \textbf{virtual water}. It takes roughly 1800 litres to grow a kilogram of wheat, and around 15\,000 litres for a kilogram of beef, mostly as rain and irrigation water for the grass and grain the cattle eat. Adding up the tap water and the virtual water gives a person's \textbf{water footprint}.

Most of a typical Australian's water footprint is in food, not in the shower. So, although shorter showers matter, especially in a drought, what we eat and what we waste make an even bigger difference. Foods grown with rain instead of irrigation, and food that is eaten instead of thrown out, both shrink the footprint.
\end{readingbox}

\rcreset
\rcq{What did Target 155 measure, and what did it leave out?}{2.5cm}
\rcq{What is virtual water? Give an example.}{2.5cm}
\rcq{Why does beef have such a large water footprint?}{2.5cm}
\rcq{Why does wasting food waste water?}{2.5cm}

\lessonpart{B}{Theory and examples}{Blue}\label{theory:3}

\subsection{Measuring a Household's Water}\label{sub:a3.1}

\begin{skillbox}{Water used}
\[ \text{water used (L)} = \text{flow rate (L/min)} \times \text{time (min)} \]
Measure a flow rate by timing how long it takes to fill a bucket of known volume.
\end{skillbox}

\begin{workedex}{A shower}
A shower fills a 9\,L bucket in 60 seconds, so its flow rate is 9\,L/min. An 8-minute shower uses $9 \times 8 = \mathbf{72\ L}$. With a water-saving shower head at 6\,L/min, the same shower uses 48\,L, a saving of 24\,L.
\end{workedex}

\begin{wedo}{One person's day}
Estimate one person's tap water for a day: a 6-minute shower at 9\,L/min; four toilet flushes at 4.5\,L; teeth brushed twice with the tap running for 2 minutes at 6\,L/min; a share of the washing machine, 30\,L. Is it under Target 155?
\workbox{5.5cm}
\end{wedo}

\subsection{Virtual Water in Food}\label{sub:a3.2}

\begin{workedex}{A burger patty}
A 150\,g beef patty, at about 15\,000\,L per kg:\\
$0.15 \times 15\,000 = \mathbf{2250\ L}$, about as much as fifteen days of tap water at Target 155.
\end{workedex}

\begin{wedo}{A sandwich}
Two slices of bread use about 80\,g of wheat (1800\,L per kg). Find the virtual water in the bread, and compare it with the burger patty.
\workbox{4.5cm}
\end{wedo}

\lessonpart{C}{Practice}{Amber}\label{prac:3}

\begin{keepext}
\task The desalination plant makes Melbourne's water supply more reliable but uses a lot of energy. Weigh up whether it makes Melbourne's water use more or less sustainable.
\answerbox{5cm}
\end{keepext}
\begin{keepext}
\task Fifteen people each use 155 litres of water a day. How much water do they use in a year? How many 9-litre buckets is that?
\workbox{4cm}
\end{keepext}


\lessonpart{D}{Challenge}{Violet}\label{ext:3}

\task A family of four each use about 200\,L of tap water a day. They want to reach Target 155.
\pt{a} How much less water must the whole family use each day? And each year?
\workbox{4cm}
\pt{b} Plan three changes that together save at least that much. Use numbers for each change.
\answerbox{6cm}

\clearpage

% ══════════════════════════════════════════════════════════════════════════════
%  TOPIC 4 — ENERGY EFFICIENCY AND THE GRID                  (deepens Lesson 5)
% ══════════════════════════════════════════════════════════════════════════════
\setcounter{section}{3}\renewcommand{\lessonart}{5}
\sectiondiv{Energy Efficiency and the Grid}{Efficiency \textbullet{} Kilowatt-hours and cost \textbullet{} Matching supply to demand \textbullet{} Storing energy}\label{sec:4}
\lessoncodes{Deepens Lesson 5 \textbullet{} VC2S10U15 \textbullet{} VC2S10U16}

\vocabhead
\begin{ctab}
\begin{tabular}{|SW{3.4cm}|W{7.0cm}|}
  \hline\hd & \thead{Word} & \thead{Definition (copy from the slide)} \\ \hline
  & \scaf{\textbf{Efficiency}\label{voc:efficiency}} & \scaf{The percentage of the energy put in that ends up as useful energy out.} \\ \hline
  & \textbf{Power}\label{voc:power} & \\ \hline
  & \textbf{Kilowatt-hour (kWh)}\label{voc:kilowatt-hour-kwh} & \\ \hline
  & \textbf{Demand}\label{voc:demand} & \\ \hline
  & \textbf{Energy storage}\label{voc:energy-storage} & \\ \hline
\end{tabular}
\end{ctab}

\lessonpart{A}{Reading and comprehension}{Teal}\label{read:4}

\begin{readingbox}{Reading --- Keeping the Lights On}
Energy is never destroyed, but it is often wasted. An old-style light globe turns most of its electrical energy into heat, not light. A car engine turns only about a quarter of the energy in its fuel into motion; the rest heats the engine, the exhaust and the air. The \textbf{efficiency} of a device is the share of the energy put in that comes out in the form we want.

\textbf{Power} is how fast energy is used, measured in watts (W) or kilowatts (kW). Electricity companies charge for energy in \textbf{kilowatt-hours} (kWh): one kilowatt of power used for one hour.

The electricity grid has a tricky job: the electricity made must match the electricity used, every second. \textbf{Demand} in Victoria peaks on summer evenings, when air conditioners and ovens are on, just as solar panels stop producing. That is why \textbf{energy storage} matters: batteries, and pumped hydro schemes such as Snowy 2.0, store spare energy from sunny, windy times and release it when it is needed.
\end{readingbox}

\rcreset
\rcq{Where does most of the energy go in an old-style light globe and in a car engine?}{2.5cm}
\rcq{What is the difference between power and energy?}{2.5cm}
\rcq{Why is a summer evening a hard time for a grid with a lot of solar power?}{2.5cm}
\rcq{How does energy storage help?}{2.5cm}

\lessonpart{B}{Theory and examples}{Blue}\label{theory:4}

\subsection{Efficiency}\label{sub:a4.1}

\begin{skillbox}{Efficiency}
\[ \text{efficiency} = \frac{\text{useful energy out}}{\text{total energy in}} \times 100\% \]
No device is 100\% efficient: some energy always ends up as heat, sound or other forms we do not want.
\end{skillbox}

\begin{workedex}{A kettle}
A kettle uses 200\,000\,J of electrical energy, and 180\,000\,J ends up heating the water.

\smallskip
Efficiency $= 180\,000 \div 200\,000 \times 100\% = \mathbf{90\%}$. The other 10\% heats the kettle and the air around it.
\end{workedex}

\begin{wedo}{A car engine}
A car burns fuel containing 400\,kJ of energy, and 100\,kJ becomes motion. Find the efficiency, and name two places the rest of the energy goes.
\workbox{4.5cm}
\end{wedo}

\subsection{Kilowatt-hours and Cost}\label{sub:a4.2}

\begin{skillbox}{Energy and cost}
\[ \text{energy (kWh)} = \text{power (kW)} \times \text{time (h)} \qquad \text{cost} = \text{energy (kWh)} \times \text{price per kWh} \]
1\,kW $=$ 1000\,W.
\end{skillbox}

\begin{workedex}{A heater}
A 2\,kW heater runs for 3 hours. Energy $= 2 \times 3 = 6$\,kWh. At 30 cents per kWh, it costs $6 \times 30 = 180$ cents, or \textbf{\$1.80}.
\end{workedex}

\begin{wedo}{Two light globes}
An old globe uses 60\,W and an LED gives the same light using 9\,W. Each is on 5 hours a day for a year (365 days). Find the energy each uses in kWh, and how much the LED saves at 30 cents per kWh.
\workbox{6cm}
\end{wedo}

\lessonpart{C}{Practice}{Amber}\label{prac:4}

\begin{keepext}
\task Pumped hydro stores energy by pumping water uphill when there is spare electricity, then letting it fall through turbines later. Explain how this helps a grid that uses a lot of solar and wind power.
\answerbox{5cm}
\end{keepext}
\begin{keepext}
\task Some people say every source of electricity has an environmental cost. Choose a renewable source and describe its environmental cost.
\answerbox{5cm}
\end{keepext}


\lessonpart{D}{Challenge}{Violet}\label{ext:4}

\task A small town wants all its electricity from rooftop solar. Its solar panels make more than the town needs from 10\,am to 3\,pm, but nothing after sunset, and its demand is highest from 6\,pm to 9\,pm. Explain the problem, propose a solution using energy storage, and describe one advantage and one disadvantage of your solution.
\answerbox{9cm}

\clearpage

\clearpage\markboth{Glossary}{Glossary}\hypertarget{glossary}{}\pdfbookmark[1]{Glossary}{glossarybm}
{\headingfont\bfseries\LARGE Glossary}\label{sec:glossary}\par\vspace{1mm}\noindent\rule{\linewidth}{0.6pt}\par
{\small Every key word in this booklet and the page it is introduced on.}\par\smallskip
\renewcommand{\arraystretch}{1.3}
\begin{longtable}{|>{\RaggedRight}p{34mm}|p{9mm}|p{\dimexpr\linewidth-43mm-6\tabcolsep-4\arrayrulewidth-1pt\relax}|}
\hline \rule{0pt}{4mm}\headingfont\bfseries Word & \headingfont\bfseries Page & \headingfont\bfseries In my own words \\ \hline \endhead
  Circular economy & \pageref{voc:circular-economy} & \\ \hline
  Demand & \pageref{voc:demand} & \\ \hline
  E-waste & \pageref{voc:e-waste} & \\ \hline
  Efficiency & \pageref{voc:efficiency} & \\ \hline
  Embodied energy & \pageref{voc:embodied-energy} & \\ \hline
  Energy storage & \pageref{voc:energy-storage} & \\ \hline
  Flow rate & \pageref{voc:flow-rate} & \\ \hline
  Flux & \pageref{voc:flux} & \\ \hline
  Groundwater & \pageref{voc:groundwater} & \\ \hline
  Kilowatt-hour (kWh) & \pageref{voc:kilowatt-hour-kwh} & \\ \hline
  Life cycle & \pageref{voc:life-cycle} & \\ \hline
  Life cycle assessment & \pageref{voc:life-cycle-assessment} & \\ \hline
  Per capita & \pageref{voc:per-capita} & \\ \hline
  Power & \pageref{voc:power} & \\ \hline
  Residence time & \pageref{voc:residence-time} & \\ \hline
  Store & \pageref{voc:store} & \\ \hline
  Target & \pageref{voc:target} & \\ \hline
  Transpiration & \pageref{voc:transpiration} & \\ \hline
  Virtual water & \pageref{voc:virtual-water} & \\ \hline
  Water footprint & \pageref{voc:water-footprint} & \\ \hline
\end{longtable}
\padtofour
\end{document}
